\documentclass[11pt,a4paper,openany]{book}

\makeatletter\def\input@path{{../_shared/}}\makeatother
\usepackage{course}

% Two boxes of this course only (kept here so the other courses do not change):
% goatbox, the idea of a section told as simply as possible, and recapbox, the big
% end-of-chapter recap to read the night before the exam (breakable: it is long).
\definecolor{goatColor}{HTML}{8A5A1F}
\definecolor{recapColor}{HTML}{5B3A8E}
\newtcolorbox{goatbox}[1][]{
  enhanced, colback=goatColor!6, colframe=goatColor,
  boxrule=0.5pt, arc=3pt, left=5pt, right=5pt, top=4pt, bottom=4pt,
  fonttitle=\bfseries\small, coltitle=white,
  colbacktitle=goatColor,
  attach boxed title to top left={yshift=-3pt, xshift=6pt},
  boxed title style={size=small, colback=goatColor, boxrule=0pt, arc=2pt},
  title=Explained to a goat, #1
}
\newtcolorbox{recapbox}[1][]{
  enhanced, breakable, colback=recapColor!5, colframe=recapColor,
  boxrule=1.2pt, arc=4pt, left=8pt, right=8pt, top=8pt, bottom=8pt,
  fonttitle=\bfseries\large, coltitle=white,
  colbacktitle=recapColor,
  title=Exam-ready recap, #1
}

\begin{document}
\raggedbottom

\coursetitle{Responsible Software}{BA5}{2026--2027}

\tableofcontents
\clearpage

These notes cover CS-290 \emph{Responsible Software} by C\'ecile Hardebolle (EPFL, Fall 2026):
how technical decisions in software have ethical consequences, and practical strategies to
anticipate and reduce them. They follow the chapters of the course (Introduction, Safety 1 \& 2,
Fairness 1 \& 2): the theory of the videos, the strategies (with their cheat sheets), and the
review questions and case studies of the Monday sessions.

Each section ends with a brown box, \emph{Explained to a goat}, that says the idea in the
simplest possible words. Each chapter ends with a big purple box, the \emph{Exam-ready recap}:
read only these boxes and you have the essentials.


% =========================================================
\chapter{Getting started: the course and ethical debt}
% =========================================================

\begin{rembox}[title={In a nutshell}]
\textbf{What is this session about?} How the course works (flipped class, notebooks, case
studies, grading, GenAI policy) and a first look at why software and AI systems cause harm:
the AI incidents monitor and the notion of \textbf{ethical debt}. \\
\textbf{Why is it useful?} Ethical debt and its causes (``veiled'' ethical agency) are the
reason the whole course exists: engineers care about ethics, but do not always see the ethical
side of their technical decisions.
\end{rembox}


\section{Course organisation}

\begin{rembox}[title={Course format: a flipped class}]
\begin{itemize}
  \item \textbf{Theory at home} (videos of theory and of \emph{strategies}, with quizzes,
        about 5 videos per week), \textbf{practice in class} so that you can get support.
        The ``lecture'' is not recorded.
  \item \textbf{Tuesday 10:15--12:00}: programming exercise (Jupyter notebook on Noto, with
        assistants). Solution provided.
  \item \textbf{Monday 15:15--17:00}: interactive review questions (ttpoll) on the videos, then
        case studies together. Solution provided.
  \item Learning loop: \textbf{Experience} (a problem to solve with programming),
        \textbf{Knowledge} (ethical and technical theory + strategies), \textbf{Habits}
        (applying strategies to cases).
\end{itemize}
\end{rembox}

\begin{rembox}[title={Contents of the course}]
\textbf{Software-related ethical challenges}:
\begin{itemize}
  \item \textbf{Safety}: harmful impacts, safety \& security, micro / meso / macro levels;
  \item \textbf{Fairness}: sources and impacts of biases, fairness measures;
  \item \textbf{Sustainability}: environmental impact, life cycle assessment for software;
  \item \textbf{Empowerment}: user agency, influence, transparency \& explainability;
  \item \textbf{Privacy}: data traces, data protection, re-identification, surveillance.
\end{itemize}
\textbf{Responsible design strategies}: stakeholder analysis, impact anticipation, values
elicitation, risk assessment, decision-making. \\
Part 1 of the semester is dense, part 2 lighter (fewer videos, more practice of strategies).
\end{rembox}

\end{document}
